\documentclass[tikz,border=4pt]{standalone}
% Paleta y fuentes compartidas por el libro y las figuras (única fuente de verdad).
\RequirePackage{xcolor}
\definecolor{oxido}{HTML}{B7410E}
\definecolor{cobre}{HTML}{D27D2D}
\definecolor{carbon}{HTML}{2B2B2B}
\definecolor{azul}{HTML}{1F5F8B}
\definecolor{arena}{HTML}{F4EDE4}

\usepackage{fontspec}
\setsansfont{FiraSans}[Extension=.otf,UprightFont=*-Regular,BoldFont=*-SemiBold,ItalicFont=*-Italic,BoldItalicFont=*-SemiBoldItalic]
\setmonofont{FiraMono}[Extension=.otf,UprightFont=*-Regular,BoldFont=*-Medium,Scale=0.95]
\usetikzlibrary{positioning,arrows.meta,calc,decorations.pathreplacing}
\begin{document}
\sffamily\small
\begin{tikzpicture}[>={Stealth[length=2.2mm]},
  cod/.style={anchor=west,font=\ttfamily\small,inner sep=1pt},
  nota/.style={anchor=west,font=\sffamily\footnotesize,align=left},
  marca/.style={rounded corners=1.5pt,line width=0.9pt}]
  % Bloque de código
  \fill[fondocodigo] (-0.25,0.45) rectangle (5.6,-3.45);
  \node[cod] at (0,0)  {let y = \{};
  \node[cod] (s) at (0.8,-1) {let x = 3;};
  \node[cod] (e) at (0.8,-2) {x + 1};
  \node[cod] at (0,-3) {\};};
  % Marcas
  \draw[marca,carbon!60,dashed] ([xshift=-2pt,yshift=2pt]s.north west) rectangle ([xshift=2pt,yshift=-2pt]s.south east);
  \draw[marca,oxido] ([xshift=-2pt,yshift=2pt]e.north west) rectangle ([xshift=2pt,yshift=-2pt]e.south east);
  % Notas
  \draw[->,thick,carbon!60] (3.4,-1) -- (6.3,-1);
  \node[nota,text=carbon!80] at (6.4,-1) {\textbf{sentencia}: termina en \texttt{;}\\y no produce valor};
  \draw[->,thick,oxido] (2.0,-2) -- (6.3,-2);
  \node[nota,text=oxido] at (6.4,-2) {\textbf{expresión final}, sin \texttt{;}:\\su valor (4) es el del bloque};
  \draw[->,thick,azul] (0.75,-3) -- (6.3,-3);
  \node[nota,text=azul] at (6.4,-3) {\textbf{el bloque \texttt{\{\ldots\}} es una expresión}:\\\texttt{y} recibe 4};
  % Contraste
  \node[nota,text=carbon!80] at (-0.25,-4.2) {Con \texttt{x + 1;} la última línea sería una sentencia y el bloque valdría \texttt{()}, el tipo unidad.};
\end{tikzpicture}
\end{document}
